Tato kapitola předkládá empirickou analýzu současného stavu sociálního dialogu v~\acsloc{ČR} s~přímým ukotvením v~širším evropském srovnání. Český trh práce se na první pohled vyznačuje velmi příznivými makroekonomickými výsledky -- vysokou mírou zaměstnanosti a~tradičně nízkou mírou příjmové nerovnosti měřenou Giniho koeficientem. Cílem této kapitoly je však poukázat na to, že tyto standardní makroindikátory mohou zastírat strukturální slabiny, které se plně vyjevují až při hlubší analýze mzdové hladiny, odpracovaných hodin, daňového zatížení, struktury volných pracovních míst a~především dosahu kolektivního vyjednávání. Právě roztříštěnost zaměstnanců a~nízké pokrytí kolektivními smlouvami tvoří jádro modelu založeného na dlouhé pracovní době a~relativně levné práci, což limituje konvergenci k~nejvyspělejším evropským ekonomikám.\par

Struktura české pracovní síly na počátku roku 2025 (tab.~\ref{tab:stav_struktura_prac_sily_cz}) demonstruje, že dominantní část tvoří zaměstnanci v~podnikatelské (mzdové) sféře, přičemž veřejný (platový) sektor je zhruba stejně velký jako podíl živnostníků. Tento specifický rys, spojený s~narůstajícím odhadovaným rozsahem tzv.~švarcsystému -- závislé ekonomické činnosti konané mimo pracovněprávní vztahy, hraje významnou roli v~oslabování vyjednávací pozice zaměstnanců, neboť vyčleňuje značnou část ekonomicky závislých pracovníků, včetně vysokopříjmových, z~ochrany zákoníkem práce a~z~dosahu kolektivních smluv.\par

\input{texparts/python/stav_struktura_prac_sily}

Další části kapitoly rozebírají diferenciaci mezd a~příjmů podle pohlaví, regionů a~odvětví, úroveň aktivní politiky zaměstnanosti a~výhledové tlaky spojené s~demografií, přeshraniční mobilitou, vzděláním, automatizací a~imigrací. Na to navazuje samostatná analýza pokrytí kolektivními smlouvami, odborové organizovanosti a~rovnováhy sil mezi zaměstnanci, zaměstnavateli a~státem. Závěr kapitoly tvoří korelační analýza a~případové studie.\par